\documentclass[12pt,a4paper]{article}

\usepackage[utf8]{inputenc}
\usepackage[T1]{fontenc}
\usepackage[french]{babel}
\usepackage[margin=2.5cm]{geometry}
\usepackage{listings}
\usepackage{xcolor}
\usepackage{hyperref}
\usepackage{parskip}
\usepackage{booktabs}
\usepackage{fancyhdr}
\usepackage{lmodern}

\setlength{\headheight}{14.5pt}

% ── style des blocs de code ───────────────────────────────────────────────────
\definecolor{codebg}{HTML}{F5F5F5}
\definecolor{keyword}{HTML}{0000CC}
\definecolor{comment}{HTML}{007700}
\definecolor{string}{HTML}{CC0000}
\definecolor{linenum}{HTML}{888888}

\lstdefinestyle{cstyle}{
    language=C,
    backgroundcolor=\color{codebg},
    basicstyle=\ttfamily\small,
    keywordstyle=\color{keyword}\bfseries,
    commentstyle=\color{comment}\itshape,
    stringstyle=\color{string},
    numberstyle=\tiny\color{linenum},
    numbers=left,
    numbersep=8pt,
    stepnumber=1,
    breaklines=true,
    breakatwhitespace=false,
    tabsize=4,
    showstringspaces=false,
    frame=single,
    framesep=4pt,
    xleftmargin=16pt,
    xrightmargin=4pt,
    captionpos=b,
}

\lstset{style=cstyle}

% ── en-tête et pied de page ───────────────────────────────────────────────────
\pagestyle{fancy}
\fancyhf{}
\rhead{Sockets 1 --- Programmation Système}
\lhead{Talkie-Walkie TCP}
\cfoot{\thepage}

% ── document ──────────────────────────────────────────────────────────────────
\begin{document}

\begin{titlepage}
    \centering
    \vspace*{2cm}
    {\Huge\bfseries Talkie-Walkie TCP\par}
    \vspace{0.5cm}
    {\Large Programmation Système en C --- Sockets 1\par}
    \vspace{2cm}
    \rule{\linewidth}{0.4pt}\\[0.4cm]
    {\large Chat semi-duplex sur socket TCP avec échange de tours\par}
    \rule{\linewidth}{0.4pt}
    \vfill
    {\large Ilyass Bougati\par}
    \vspace{0.3cm}
    {\large \today\par}
\end{titlepage}

\tableofcontents
\newpage

% ── Introduction ──────────────────────────────────────────────────────────────
\section{Introduction}

Ce projet implémente une application de chat en mode \textbf{semi-duplex}
(talkie-walkie) sur une socket TCP. Le principe est simple : à l'image d'un
vrai talkie-walkie, un seul des deux interlocuteurs peut parler à la fois. Le
client envoie un message, attend la réponse du serveur, puis envoie à nouveau
--- et ainsi de suite. Aucun fil d'exécution supplémentaire n'est nécessaire ;
les appels à \texttt{send} et \texttt{recv} sont simplement enchaînés de
manière alternée.

Le projet se compose de deux programmes :

\begin{itemize}
    \item \textbf{server} --- se met en écoute sur un port, accepte les
          connexions entrantes une par une et prend part à la conversation.
    \item \textbf{client} --- se connecte à l'adresse et au port indiqués
          et initie la conversation.
\end{itemize}

% ── Architecture ──────────────────────────────────────────────────────────────
\section{Architecture réseau}

\subsection{Établissement de la connexion}

La connexion suit le modèle TCP classique :

\begin{enumerate}
    \item Le serveur crée une socket (\texttt{SOCK\_STREAM}), lie
          (\texttt{bind}) l'adresse \texttt{INADDR\_ANY} au port demandé,
          puis se met en écoute (\texttt{listen}).
    \item Le client crée sa propre socket et appelle \texttt{connect} vers
          l'adresse IP et le port du serveur.
    \item Le serveur appelle \texttt{accept}, qui retourne un descripteur de
          fichier dédié à cette connexion.
\end{enumerate}

L'option \texttt{SO\_REUSEADDR} est positionnée sur le serveur pour éviter
l'erreur \texttt{Address already in use} lors des redémarrages rapides.

\subsection{Échange des pseudonymes}

Avant d'entrer dans la boucle de chat, les deux parties échangent leur
pseudonyme :

\begin{enumerate}
    \item Le serveur saisit son nom et l'envoie au client.
    \item Le client reçoit ce nom, saisit le sien et l'envoie en retour.
\end{enumerate}

Cet handshake permet à chaque côté d'afficher les messages sous la forme
\texttt{<nom>: <message>}.

\subsection{Protocole semi-duplex}

La boucle de conversation est strictement alternée :

\begin{center}
\begin{tabular}{@{}lll@{}}
\toprule
Tour & Côté client & Côté serveur \\
\midrule
1 & \texttt{send(message)} & \texttt{recv(message)} \\
2 & \texttt{recv(réponse)} & \texttt{send(réponse)} \\
3 & \texttt{send(message)} & \texttt{recv(message)} \\
\ldots & \ldots & \ldots \\
\bottomrule
\end{tabular}
\end{center}

La commande \texttt{/quit} met fin à la conversation des deux côtés.

% ── Signaux ───────────────────────────────────────────────────────────────────
\section{Gestion des signaux}

Le serveur installe un gestionnaire pour \texttt{SIGINT} (Ctrl+C) via
\texttt{signal(SIGINT, handle\_sigint)}. Ce gestionnaire ferme proprement
le descripteur d'écoute avant de quitter, évitant ainsi que le port reste
occupé.

% ── Utilisation ───────────────────────────────────────────────────────────────
\section{Utilisation}

\begin{lstlisting}[language=bash, style=cstyle]
# Terminal 1 : lancer le serveur sur le port 4242
./server 4242

# Terminal 2 : connecter le client
./client 127.0.0.1 4242
\end{lstlisting}

% ── Code source ───────────────────────────────────────────────────────────────
\section{Code source}

\subsection{\texttt{server.c}}

\begin{lstlisting}[caption={\texttt{server.c} --- serveur de chat semi-duplex}]
#include <stdio.h>
#include <stdlib.h>
#include <string.h>
#include <unistd.h>
#include <signal.h>
#include <arpa/inet.h>

#define BUF_SIZE 1024
#define NAME_SIZE 64

static int server_fd = -1;

static void handle_sigint(int sig) {
    (void)sig;
    printf("\nClosing the server application...\n");
    if (server_fd != -1)
        close(server_fd);
    exit(0);
}

void chat(int client_fd) {
    char my_name[NAME_SIZE];
    char their_name[NAME_SIZE];
    char buf[BUF_SIZE];
    ssize_t n;

    printf("Enter your username: ");
    fflush(stdout);
    if (!fgets(my_name, sizeof(my_name), stdin))
        return;
    my_name[strcspn(my_name, "\n")] = '\0';

    send(client_fd, my_name, strlen(my_name), 0);

    n = recv(client_fd, their_name, sizeof(their_name) - 1, 0);
    if (n <= 0)
        return;
    their_name[n] = '\0';

    printf("%s joined the chat. They go first.\n", their_name);

    while (1) {
        n = recv(client_fd, buf, sizeof(buf) - 1, 0);
        if (n <= 0) {
            printf("%s disconnected.\n", their_name);
            break;
        }
        buf[n] = '\0';

        if (strcmp(buf, "/quit") == 0) {
            printf("%s left the chat.\n", their_name);
            break;
        }
        printf("%s: %s\n", their_name, buf);

        printf("%s: ", my_name);
        fflush(stdout);
        if (!fgets(buf, sizeof(buf), stdin))
            break;
        buf[strcspn(buf, "\n")] = '\0';

        send(client_fd, buf, strlen(buf), 0);

        if (strcmp(buf, "/quit") == 0)
            break;
    }

    close(client_fd);
    printf("Chat ended. Waiting for a new connection...\n\n");
}

int main(int argc, char *argv[]) {
    if (argc != 2) {
        fprintf(stderr, "Usage: %s <port>\n", argv[0]);
        exit(EXIT_FAILURE);
    }

    int port = atoi(argv[1]);
    if (port <= 0 || port > 65535) {
        fprintf(stderr, "Invalid port: %s\n", argv[1]);
        exit(EXIT_FAILURE);
    }

    signal(SIGINT, handle_sigint);

    int client_fd;
    struct sockaddr_in address;
    socklen_t addr_len = sizeof(address);

    server_fd = socket(AF_INET, SOCK_STREAM, 0);
    if (server_fd < 0) {
        perror("socket");
        exit(EXIT_FAILURE);
    }

    int opt = 1;
    setsockopt(server_fd, SOL_SOCKET, SO_REUSEADDR, &opt, sizeof(opt));

    address.sin_family = AF_INET;
    address.sin_addr.s_addr = INADDR_ANY;
    address.sin_port = htons(port);

    if (bind(server_fd, (struct sockaddr *)&address, sizeof(address)) < 0) {
        perror("bind");
        exit(EXIT_FAILURE);
    }

    if (listen(server_fd, 1) < 0) {
        perror("listen");
        exit(EXIT_FAILURE);
    }

    printf("Server listening on port %d...\n\n", port);

    while (1) {
        client_fd = accept(server_fd, (struct sockaddr *)&address, &addr_len);
        if (client_fd < 0) {
            perror("accept");
            continue;
        }
        printf("Client connected from %s\n", inet_ntoa(address.sin_addr));
        chat(client_fd);
    }

    close(server_fd);
    return 0;
}
\end{lstlisting}

\subsection{\texttt{client.c}}

\begin{lstlisting}[caption={\texttt{client.c} --- client de chat semi-duplex}]
#include <stdio.h>
#include <stdlib.h>
#include <string.h>
#include <unistd.h>
#include <arpa/inet.h>

#define BUF_SIZE 1024
#define NAME_SIZE 64

void chatc(int sock_fd) {
    char my_name[NAME_SIZE];
    char their_name[NAME_SIZE];
    char buf[BUF_SIZE];
    ssize_t n;

    n = recv(sock_fd, their_name, sizeof(their_name) - 1, 0);
    if (n <= 0) {
        fprintf(stderr, "Server closed connection.\n");
        return;
    }
    their_name[n] = '\0';

    printf("Enter your username: ");
    fflush(stdout);
    if (!fgets(my_name, sizeof(my_name), stdin))
        return;
    my_name[strcspn(my_name, "\n")] = '\0';

    send(sock_fd, my_name, strlen(my_name), 0);

    printf("Connected to %s. You go first!\n\n", their_name);

    while (1) {
        printf("%s: ", my_name);
        fflush(stdout);
        if (!fgets(buf, sizeof(buf), stdin))
            break;
        buf[strcspn(buf, "\n")] = '\0';

        send(sock_fd, buf, strlen(buf), 0);

        if (strcmp(buf, "/quit") == 0)
            break;

        n = recv(sock_fd, buf, sizeof(buf) - 1, 0);
        if (n <= 0) {
            printf("%s disconnected.\n", their_name);
            break;
        }
        buf[n] = '\0';

        if (strcmp(buf, "/quit") == 0) {
            printf("%s left the chat.\n", their_name);
            break;
        }
        printf("%s: %s\n", their_name, buf);
    }
}

int main(int argc, char *argv[]) {
    if (argc != 3) {
        fprintf(stderr, "Usage: %s <host> <port>\n", argv[0]);
        exit(EXIT_FAILURE);
    }

    int port = atoi(argv[2]);
    if (port <= 0 || port > 65535) {
        fprintf(stderr, "Invalid port: %s\n", argv[2]);
        exit(EXIT_FAILURE);
    }

    int sock_fd;
    struct sockaddr_in server_addr;

    sock_fd = socket(AF_INET, SOCK_STREAM, 0);
    if (sock_fd < 0) {
        perror("socket");
        exit(EXIT_FAILURE);
    }

    server_addr.sin_family = AF_INET;
    server_addr.sin_port = htons(port);
    if (inet_pton(AF_INET, argv[1], &server_addr.sin_addr) <= 0) {
        fprintf(stderr, "Invalid address: %s\n", argv[1]);
        exit(EXIT_FAILURE);
    }

    if (connect(sock_fd, (struct sockaddr *)&server_addr,
                sizeof(server_addr)) < 0) {
        perror("connect");
        exit(EXIT_FAILURE);
    }

    chatc(sock_fd);

    close(sock_fd);
    return 0;
}
\end{lstlisting}

\subsection{Makefile}

\begin{lstlisting}[language=make, caption={Makefile --- compilation des deux binaires}]
CC = gcc
CFLAGS = -Wall -Wextra

all: server client

server: server.c
	$(CC) $(CFLAGS) -o $@ $<

client: client.c
	$(CC) $(CFLAGS) -o $@ $<

clean:
	rm -f server client

.PHONY: all clean
\end{lstlisting}

% ── Bilan ─────────────────────────────────────────────────────────────────────
\section{Bilan}

Ce projet illustre les bases de la programmation réseau en C : création de
socket, liaison, écoute, connexion et échange de données. Le modèle
semi-duplex est intentionnellement simple --- aucun thread, aucun
multiplexage --- ce qui permet de se concentrer sur le cycle de vie d'une
connexion TCP et sur le protocole applicatif minimal (échange de noms,
boucle alternée, commande \texttt{/quit}).

La limitation principale est précisément ce caractère bloquant : si l'un
des deux interlocuteurs tarde à répondre, l'autre est figé dans son appel
\texttt{recv}. C'est cette contrainte que lève le projet suivant
(\emph{Sockets 2}) en introduisant des fils d'exécution dédiés à la
réception.

\end{document}
